% !TeX root = ../navier-stokes-manuscript.tex
\subsection{Vortex Stretching}\label{sub:ThePerfectlyStretchingVortex}

The momentum equation puts four pieces into one evolution:
\[
  \underbrace{\partial_tu}_{\text{local acceleration}}
  +
  \underbrace{(u\cdot\nabla)u}_{\text{nonlinear transport}}
  =
  \underbrace{-\nabla p}_{\text{pressure}}
  +
  \underbrace{\nu\Delta u}_{\text{viscous diffusion}}.
\]
For the regularity problem, the central competition is between nonlinear transport and viscosity.
The hope is that viscosity supplies all the smoothing needed to keep the velocity smooth: it must smooth the fluid at least as quickly as nonlinear transport sharpens it.

In three dimensions, the difficulty comes from vortex stretching.
When a region of concentrated rotation is stretched along its axis, incompressibility reduces its transverse cross-sectional area.
If a nontrivial transverse velocity difference persists while that distance shrinks, the corresponding spatial gradient can grow.
Finite kinetic energy does not exclude that growth because it controls total squared velocity, not the distance across which the velocity changes.
A finite-time singularity would require the sharpening to continue until some spatial derivative becomes unbounded.

The stretching mechanism appears when the momentum equation is curled.
Define
\[
  \omega:=\nabla\times u,
  \qquad
  S:=\frac12\bigl(\nabla u+(\nabla u)^{\mathsf T}\bigr).
\]
Taking the curl of Navier--Stokes gives
\begin{equation}\label{eq:VorticityEquation}
  \partial_t\omega+(u\cdot\nabla)\omega
  =(\omega\cdot\nabla)u+\nu\Delta\omega
  =S\omega+\nu\Delta\omega.
\end{equation}
The left side follows the vorticity carried by the moving fluid.
The term \(S\omega\) is vortex stretching.
When the vorticity is aligned with an extensional direction of the strain, \(S\omega\) amplifies it.
Taking the scalar product with \(\omega\) gives
\[
  \frac12(\partial_t+u\cdot\nabla)|\omega|^2
  =\omega\cdot S\omega
  +\nu\,\omega\cdot\Delta\omega.
\]
The first term can increase the strength of the vorticity.
The second is the viscous response to its spatial variation.
The original competition has reappeared in the quantity that measures concentrated rotation.

To see what happens as the distance across the vortex shrinks, rescale the full equation.
For \(\lambda>1\), set
\[
  u_\lambda(x,t)=\lambda u(\lambda x,\lambda^2t),
  \qquad
  p_\lambda(x,t)=\lambda^2p(\lambda x,\lambda^2t).
\]
Every piece of the momentum equation acquires the same factor:
\[
\begin{aligned}
  \partial_tu_\lambda(x,t)
  &=\lambda^3(\partial_tu)(\lambda x,\lambda^2t),
  &
  (u_\lambda\cdot\nabla)u_\lambda(x,t)
  &=\lambda^3\bigl((u\cdot\nabla)u\bigr)(\lambda x,\lambda^2t),
  \\
  \nabla p_\lambda(x,t)
  &=\lambda^3(\nabla p)(\lambda x,\lambda^2t),
  &
  \Delta u_\lambda(x,t)
  &=\lambda^3(\Delta u)(\lambda x,\lambda^2t).
\end{aligned}
\]
The transformation preserves Navier--Stokes and produces a solution with rescaled initial data.
It compares solutions rather than asserting that \(u_\lambda\) is a later state of the original history.
The factor \(\lambda>1\) displays the velocity at spatial scale \(\lambda^{-1}\), time scale \(\lambda^{-2}\), and velocity scale \(\lambda\).

Write \(U\) for the velocity amplitude and \(\ell\) for the active spatial width.
In the critical-scaling spike picture, \(\lambda\simeq\ell^{-1}\), so
\[
  U\sim\ell^{-1},
  \qquad
  |\nabla u|\sim\ell^{-2},
  \qquad
  \tau\sim\ell^2.
\]
If the spike occupies a region of volume comparable to \(\ell^3\), the kinetic energy carried by that packet scales as
\[
  E_\ell\sim U^2\ell^3\sim\ell\longrightarrow0.
\]
Its velocity amplitude and gradients can diverge even as the energy of the shrinking packet tends to zero.
This is the critical spike picture supplied by the scaling, not a conclusion about every concentrating velocity history.

This scaling identifies the natural level at which concentration should be measured.
For a homogeneous Sobolev norm,
\[
  \norm{u_\lambda(t)}_{\dot H^s}^2
  =\lambda^{2s-1}
  \norm{u(\lambda^2t)}_{\dot H^s}^2.
\]
The exponent vanishes when
\[
  2s-1=0,
  \qquad
  s=\frac12.
\]
With \(\Lambda:=(-\Delta)^{1/2}\), define
\begin{equation}\label{eq:BodyCriticalHeightDefinition}
  H(t)
  :=\frac12\norm{u(t)}_{\dot H^{1/2}}^2
  =\frac12\norm{\Lambda^{1/2}u(t)}_2^2.
\end{equation}
We call \(H(t)\) the critical height.

Meanwhile, the basic energy identity remains
\begin{equation}\label{eq:BodyEnergyIdentity}
  \frac12\norm{u(t)}_2^2
  +\nu\int_0^t\norm{\nabla u(s)}_2^2\,ds
  =\frac12\norm{u_0}_2^2.
\end{equation}
Under the same rescaling,
\[
  \norm{u_\lambda(t)}_2^2
  =\lambda^{-1}\norm{u(\lambda^2t)}_2^2,
  \qquad
  \norm{\nabla u_\lambda(t)}_2^2
  =\lambda\norm{\nabla u(\lambda^2t)}_2^2,
\]
while
\[
  H_\lambda(t)=H(\lambda^2t).
\]
A smaller-scale profile can carry less kinetic energy, a larger first-derivative norm, and the same critical height.
This is why finite kinetic energy does not exclude the sharpening that vortex stretching can produce.

The critical spike picture and the general critical-height route must be kept separate.
The critical-height route tracks growth of \(H\), but that growth alone does not imply \(\norm{u}_{L^\infty}\to\infty\).
A velocity field can develop increasingly steep spatial structure while its pointwise amplitude remains bounded.
Velocity-amplitude blow-up follows from the critical spike law \(U\sim\ell^{-1}\); outside that model it requires additional information.

The competition can also be read at the critical height itself.
Set
\[
  \mathcal P_{1/2}(t)
  :=-
  \operatorname{Re}
  \left\langle
    \Lambda^{1/2}\bigl((u\cdot\nabla)u+\nabla p\bigr),
    \Lambda^{1/2}u
  \right\rangle_{L^2}.
\]
Applying \(\Lambda^{1/2}\) to Navier--Stokes and pairing with \(\Lambda^{1/2}u\) gives
\begin{equation}\label{eq:BodyCriticalHeightBalance}
  H'(t)
  +\nu\norm{\Lambda^{3/2}u(t)}_2^2
  =\mathcal P_{1/2}(t).
\end{equation}
The production term can raise the critical height, while viscosity dissipates it through one additional spatial derivative.
The same scaling that identifies this height also fixes the time on which viscosity responds to each shrinking spatial scale.
